\documentclass[tikz,border=12pt]{standalone}
\usepackage{fontspec}
\usepackage[mathrm=sym]{unicode-math}
\setmainfont{texgyrepagella-regular.otf}[
    BoldFont = texgyrepagella-bold.otf,
    ItalicFont = texgyrepagella-italic.otf,
    BoldItalicFont = texgyrepagella-bolditalic.otf
]
\setmathfont{texgyrepagella-math.otf}[AutoFakeBold = 2.5]
\usepackage{tikz}
\usetikzlibrary{arrows.meta, calc}

\begin{document}
\begin{tikzpicture}[>=Stealth, line cap=round, line join=round]

  % Scale factors: 1 unit of alpha = 1.1 cm
  % Origin of coordinates at (0, 0)
  
  % Horizontal axis (alpha)
  \draw[line width=0.85pt, ->] (-5.6, 0) -- (6.0, 0) node[below=3pt, font=\large] {$\alpha$};

  % Vertical axis (epsilon_i) on left
  \draw[line width=0.85pt, ->] (-5.6, 0) -- (-5.6, 8.5) node[below left=2pt, font=\large] {$\epsilon_i$};

  % Center vertical line at alpha = 0
  \draw[line width=0.85pt] (0, 0) -- (0, 8.5);

  % Alpha ticks and labels with explicit braces
  \foreach \a/\lbl in {
    -4.5/{-9/2}, -3.5/{-7/2}, -2.5/{-5/2}, -1.5/{-3/2}, -0.5/{-1/2},
     0.5/{1/2},   1.5/{3/2},   2.5/{5/2},   3.5/{7/2},   4.5/{9/2}
  } {
    \draw[line width=0.6pt] (\a, -0.08) -- (\a, 0.08);
    \node[below=3pt, font=\footnotesize] at (\a, 0) {$\lbl$};
  }

  % 1. Top parabola: j = 5/2
  \draw[line width=1.0pt, domain=-2.6:2.6, samples=50, smooth]
    plot (\x, {7.4 - 0.10*\x*\x});
  \node[left=4pt, font=\normalsize] at (-1.8, 7.3) {$j = 5/2$};
  \foreach \x in {-2.5, -1.5, -0.5, 0.5, 1.5, 2.5} {
    \filldraw[fill=white, draw=black, line width=0.8pt] (\x, {7.4 - 0.10*\x*\x}) circle (2.2pt);
  }

  % 2. Second parabola: j = 9/2
  \draw[line width=1.0pt, domain=-4.6:4.6, samples=70, smooth]
    plot (\x, {6.1 - 0.09*\x*\x});
  \node[left=4pt, font=\normalsize] at (-3.4, 5.0) {$j = 9/2$};
  \foreach \x in {-4.5, -3.5, -2.5, -1.5, -0.5, 0.5, 1.5, 2.5, 3.5} {
    \filldraw[fill=white, draw=black, line width=0.8pt] (\x, {6.1 - 0.09*\x*\x}) circle (2.2pt);
  }
  % alpha = 4.5 is FILLED!
  \fill (4.5, {6.1 - 0.09*4.5*4.5}) circle (2.5pt);

  % 3. Fermi surface for I = 0 (solid horizontal line)
  \draw[line width=0.9pt] (-5.0, 3.6) -- (5.0, 3.6);
  \node[above=2pt, font=\normalsize, fill=white, inner sep=1.0pt] at (4.1, 3.6) {$I = 0$};

  % 4. Third parabola: j = 3/2
  \draw[line width=1.0pt, domain=-1.6:1.6, samples=40, smooth]
    plot (\x, {2.9 - 0.09*\x*\x});
  \node[right=6pt, font=\normalsize] at (1.6, 2.6) {$j = 3/2$};
  % alpha = -1.5 is OPEN circle
  \filldraw[fill=white, draw=black, line width=0.8pt] (-1.5, {2.9 - 0.09*1.5*1.5}) circle (2.2pt);
  % alpha = -0.5, 0.5, 1.5 are FILLED circles
  \foreach \x in {-0.5, 0.5, 1.5} {
    \fill (\x, {2.9 - 0.09*\x*\x}) circle (2.2pt);
  }

  % 5. Bottom parabola: j = 5/2
  \draw[line width=1.0pt, domain=-2.6:2.6, samples=50, smooth]
    plot (\x, {1.9 - 0.09*\x*\x});
  \node[right=6pt, font=\normalsize] at (2.4, 1.25) {$j = 5/2$};
  \foreach \x in {-2.5, -1.5, -0.5, 0.5, 1.5, 2.5} {
    \fill (\x, {1.9 - 0.09*\x*\x}) circle (2.2pt);
  }

  % 6. Tilted Fermi surface for I = 6 (dashed inclined line)
  \draw[line width=0.9pt, dashed] (-5.2, 0.8) -- (4.9, 4.45);
  \node[rotate=28, font=\normalsize, fill=white, inner sep=1.2pt] at (-3.7, 1.95) {$I = 6 = 9/2 + 3/2$};

\end{tikzpicture}
\end{document}
